\section{Especialización de dominio frente a modalidad}
\label{sec:disc_especializacion}
\pendiente{Se discute por qué el factor determinante no es que el modelo sea visión--lenguaje, sino que esté
especializado en dermatología, y cómo se ubica esa conclusión respecto de los trabajos revisados en la
Sección~\ref{sec:estado_arte}.}

% --- Material de 4.6, párrafo 1 (revisión del tutor del 26/10/2026: 4.6 y 5.6 quedan como balance breve y lo
% argumentativo pasa a este capítulo). Para integrar al redactar esta sección; no se imprime.
% Los resultados del capítulo permiten responder la primera pregunta de investigación
% (Sección~\ref{sec:intro_objetivos}): en la prueba de \emph{HAM10000}, ningún \gls{vlm} sin una cabeza entrenada sobre
% ese conjunto clasifica las lesiones tan bien como un clasificador supervisado. Los \glspl{vlm} generativos de propósito general quedan en el nivel del azar y
% colapsan hacia una clase que depende del modelo y de la estrategia, y los controles muestran que ese comportamiento no
% se explica por el desbalance de las clases ni por la modalidad de la imagen. DermLIP, en cambio, muestra que la
% especialización en dermatología aporta una representación útil, pero solo se acerca a los supervisados con una cabeza
% entrenada sobre \emph{HAM10000}. Las diferencias con los modelos sin ajuste son, además, significativas en McNemar,
% salvo la de DenseNet frente a DermLIP en sensibilidad macro, y coinciden con lo que informa la literatura.
